
\begin{frame}{Выводы работы}
	\begin{itemize}
	 \item 
	 Поставлена и решена задача принятия решения о выборе наиболее качественных объектов на основе нечётких экспертных данных об их характеристиках, сформулированная как задача на минимум возможности принять неверное 
	 решение (ошибиться);	
	\item 
	Разработан численный метод, позволяющий решить эту задачу за время порядка $O(nm)$ вместо $O\big(\abs{X}^{nm}\big)$ в случае полного перебора всех элементраных событий. %, где $n$ --- количество объектов, $m$ --- количество параметров у каждого объекта, а $\abs{X}$ --- размер множества $X$ значений параметров.
	\item
	Поставлена и решена задача нахождения коллективного мнения экспертов как задача нахождения верхней грани на множестве возможностных распределений (экспертных оценок), используя введённое на нём отношение квазипорядка.
	\item
	Предложен и обоснован численный метод решения этой задачи за время порядка $O\big(R(nm)^2\big)$ вместо $O\big(R\abs{X}^{2nm}\big)$ в случае построения совместного распределения возможностей всех параметров всех объектов. %, где $n,m,\abs{X}$~--- те же, что выше, а $R$ --- количество экспертов. 
      \end{itemize}
\end{frame} %===========================

%\begin{frame}{Выводы (задача выбора объектов)}
%	\begin{itemize}
		
%	\end{itemize}
%\end{frame} %===========================

% === final slide ====
\begin{frame}{Спасибо за внимание!}
	\begin{center}
		\includegraphics[width=0.5\linewidth]{./pic/biber_final}
	\end{center}
\end{frame}
% =============

\begin{frame}{Метод векторов предпочтений}
	\vspace{-3mm}
	{ \small Коллективное мнение экспертов с помощью векторов предпочтений аналогично случаю матриц попарных сравнений ($r = 1 \ldots R$, $J = \{1\ldots\abs{T}+1\})$:}
	\begin{gather*}
		s^{(r)}_j = \sum_{t \in T} \mathlarger{\mathlarger{\chi} }_{A_j^{(r)}}(t),
		 \mathsmaller{\text{\small причём } \p^{(r)}(t_{\abs{T}+1}) \define= 0, \; j \in J. } 
		 \\ \text{Здесь } A^{(r)}_j = \{t \in T: \p^{(r)}(t) \leq \p^{(r)}(t_j)\}. 
	\end{gather*}
	Вектор предпочтений $s$ должен удовлетворять условию:
	%\\ (1) $\max s_j= |T|$ (нормировка возможности);
	\\[1.2ex] ($\star$) $|J_i| = i$, где $J_i = \{j \in J: s_j \leq i\}$, $i \in \{s_1 \ldots s_{\abs{T}+1}\}$.
	
	\textbf{Задача}: $\displaystyle s_* = \arg \underset{s} \min \sum_{r=1}^R \rho(s - s^{(r)})$, где  $\displaystyle \rho(s, s') = \big( \sum_{j=1}^{\abs{T}+1}(s_{j} - {s'}_{j})^2 \big)^{1/2}$.
	
	\textbf{Алгоритм решения}:  берём  вектор $ \ol s =  \frac{1}{R} \sum_{r=1}^R s^{(r)}$ и удовлетворяем условию ($\star$) с помощью округления и (если требуется) дополнительного изменения его координат.
\end{frame} %===========================
